% GENERATED FILE. Edit canonical lessons, then run npm run generate:books.
% canonical-source-hash: fnv1a64:fae51af15ebb9084
% canonical-lessons: HI-C47-teacher, HI-C47-student, HI-C47-farmer, HI-C47-guest, HI-C47-neighbour

\chapter{Who Someone Is}
\label{ch:hi-roles}
\hlchaptermodality{47}

\begin{chapteropening}
\textbf{By the end of this chapter:} \emph{I can say what someone's work is in Hindi, and tell an inherited word for a person from a borrowed one.}

Complete HI-C47-neighbour and say all five words in order.
\end{chapteropening}

\section[śikṣak]{\hi{शिक्षक} (śikṣak) — a teacher}

\label{lesson:HI-C47-teacher}



\begin{quote}
Before the new run of words: what did \emph{nāk} mean, and what did \emph{galā} mean?
\end{quote}

\subsection*{You'll want to know: \hi{शिक्षक}}
\textbf{\hi{शिक्षक}} (\emph{śikṣak}) — \textquotedblleft{}a teacher\textquotedblright{}.

Sanskrit, and it comes apart cleanly in the hand: the root \emph{śikṣ-}, 'to learn, to be taught', with the agent ending \emph{-aka} — the one who brings the learning about. The same root gives \emph{śikṣā}, education, and \emph{śikṣit}, educated.

Everyday speech reaches for \emph{ustād} or \emph{guru} depending on what is being taught and by whom, but \hi{शिक्षक} is the word on the school form, and the feminine \emph{śikṣikā} is built from it by the ending alone.

\begin{grammarlens}[title={What you've built}]
The first of five words for who someone is.
\end{grammarlens}

\subsection*{Your turn}
Say these aloud:

\begin{itemize}
\raggedright
  \item \emph{śikṣak}
  \item it once more, slowly
  \item \emph{śikṣak}, then \emph{galā}, so the teacher and the voice sit together
\end{itemize}

\begin{tcolorbox}[breakable,colback=teal!4,colframe=teal!35!black,title={Before you move on}]
What does \hi{शिक्षक} mean? (\textquotedblleft{}a teacher\textquotedblright{}.) Say it once more.
\end{tcolorbox}
\section[chātra]{\hi{छात्र} (chātra) — a student}

\label{lesson:HI-C47-student}



\begin{quote}
Before the new one: what did \emph{śikṣak} mean?
\end{quote}

\subsection*{You'll want to know: \hi{छात्र}}
\textbf{\hi{छात्र}} (\emph{chātra}) — \textquotedblleft{}a student\textquotedblright{}.

Sanskrit \emph{chātra}, and the picture inside it is unexpected. \emph{Chattra} is a parasol, and the traditional Indian etymologists read \emph{chātra} as the one who stands in the teacher's shade. Whether that is the true history is argued about; that it is the story the tradition told is not.

The feminine is \emph{chātrā}. The transparent alternative is \emph{vidyārthī}, 'one who seeks knowledge', which says on its face what \hi{छात्र} only hints at.

\begin{grammarlens}[title={What you've built}]
Two: a teacher and a student.
\end{grammarlens}

\subsection*{Your turn}
Say these aloud:

\begin{itemize}
\raggedright
  \item \emph{chātra}
  \item it once more, slowly
  \item \emph{chātra}, then \emph{śikṣak}, the pair that needs each other
\end{itemize}

\begin{tcolorbox}[breakable,colback=teal!4,colframe=teal!35!black,title={Before you move on}]
What does \hi{छात्र} mean? (\textquotedblleft{}a student\textquotedblright{}.) Say it once more.
\end{tcolorbox}
\section[kisān]{\hi{किसान} (kisān) — a farmer}

\label{lesson:HI-C47-farmer}



\begin{quote}
Before the new one: what did \emph{chātra} mean?
\end{quote}

\subsection*{You'll want to know: \hi{किसान}}
\textbf{\hi{किसान}} (\emph{kisān}) — \textquotedblleft{}a farmer\textquotedblright{}.

From \emph{kṛṣāṇa}, on the Sanskrit root \emph{kṛṣ-}, 'to plough, to drag through the ground' — the same root behind \emph{kṛṣi}, agriculture. A farmer, in this word, is simply the one who drags the plough.

Latin reached for the same picture and got there with different material: \emph{sulcus}, a furrow, is not related to \emph{kṛṣ-} at all. Two families, no shared word, and the identical image — which is worth noticing precisely because it proves nothing about descent.

\begin{grammarlens}[title={What you've built}]
Three, and each one names work rather than family.
\end{grammarlens}

\subsection*{Your turn}
Say these aloud:

\begin{itemize}
\raggedright
  \item \emph{kisān}
  \item it once more, slowly
  \item \emph{kisān}, then \emph{chātra}, then \emph{śikṣak}, back down the run
\end{itemize}

\begin{tcolorbox}[breakable,colback=teal!4,colframe=teal!35!black,title={Before you move on}]
What does \hi{किसान} mean? (\textquotedblleft{}a farmer\textquotedblright{}.) Say it once more.
\end{tcolorbox}
\section[mehmān]{\hi{मेहमान} (mehmān) — a guest}

\label{lesson:HI-C47-guest}



\begin{quote}
Before the new one: what did \emph{kisān} mean?
\end{quote}

\subsection*{You'll want to know: \hi{मेहमान}}
\textbf{\hi{मेहमान}} (\emph{mehmān}) — \textquotedblleft{}a guest\textquotedblright{}.

Persian \emph{mihmān}, and Hindi took the whole vocabulary of hospitality along with it: \emph{mehmān-nawāzī}, hospitality, is literally the cherishing of guests.

The inherited word is \emph{atithi}, and it is built out of a small argument: \emph{a-tithi}, 'without a date' — a guest being, by definition, someone who has not told you when they are coming. Two words, two theories of the thing. One names the arrival, the other names the welcome.

\begin{grammarlens}[title={What you've built}]
Four, and this one arrives at your door.
\end{grammarlens}

\subsection*{Your turn}
Say these aloud:

\begin{itemize}
\raggedright
  \item \emph{mehmān}
  \item it once more, slowly
  \item \emph{mehmān}, then \emph{kisān}, so the visitor and the host sit together
\end{itemize}

\begin{tcolorbox}[breakable,colback=teal!4,colframe=teal!35!black,title={Before you move on}]
What does \hi{मेहमान} mean? (\textquotedblleft{}a guest\textquotedblright{}.) Say it once more.
\end{tcolorbox}
\section[paṛosī]{\hi{पड़ोसी} (paṛosī) — a neighbour}

\label{lesson:HI-C47-neighbour}



\begin{quote}
Before the new one: what did \emph{mehmān} mean?
\end{quote}

\subsection*{You'll want to know: \hi{पड़ोसी}}
\textbf{\hi{पड़ोसी}} (\emph{paṛosī}) — \textquotedblleft{}a neighbour\textquotedblright{}.

Built on \emph{paṛos}, the neighbourhood, with the \emph{-ī} that turns a place into a person — the very ending that makes \hi{आदमी} out of Adam.

Behind \emph{paṛos} is Sanskrit \emph{prati-veśa}, 'the dwelling opposite', on \emph{veś-}, 'to enter, to settle'. So a neighbour here is not merely the one who is near you. It is the one whose door faces yours.

\begin{grammarlens}[title={What you've built}]
Five, and the run is closed: a teacher, a student, a farmer, a guest, a neighbour.
\end{grammarlens}

\subsection*{Your turn}
Say these aloud:

\begin{itemize}
\raggedright
  \item \emph{paṛosī}
  \item it once more, slowly
  \item all five in order — \emph{śikṣak}, \emph{chātra}, \emph{kisān}, \emph{mehmān}, \emph{paṛosī}
\end{itemize}

\begin{tcolorbox}[breakable,colback=teal!4,colframe=teal!35!black,title={Before you move on}]
What does \hi{पड़ोसी} mean? (\textquotedblleft{}a neighbour\textquotedblright{}.) Say it once more.
\end{tcolorbox}
